% Created 2026-09-05 Sat 13:25
% Intended LaTeX compiler: xelatex
\documentclass[12pt]{article}
\usepackage[utf8]{inputenc}
\usepackage[T1]{fontenc}
\usepackage{graphicx}
\usepackage{longtable}
\usepackage{wrapfig}
\usepackage[margin=1in]{geometry}
\usepackage{rotating}
\usepackage{setspace}
\onehalfspacing
\usepackage{fontspec}
\setmainfont{Libertinus Serif}
\usepackage[normalem]{ulem}
\usepackage{amsmath}
\usepackage{parskip}
\usepackage{amssymb}
\usepackage{capt-of}
\setcounter{secnumdepth}{4} % Tells LaTeX to number down to paragraphs
\setcounter{tocdepth}{4}
\usepackage{hyperref}
\usepackage[style=oscola, backend=biber]{biblatex}
\addbibresource{/Users/krishnajani/Documents/Libib.bib}
\usepackage{csquotes}
\usepackage[british]{babel}
\setlength{\parskip}{0.5\baselineskip}
\author{Krishna Jani}
\date{\today}
\title{Mobilox Innovations v. Kirusa Software}

\begin{document}
\maketitle

\section*{Facts}

\begin{itemize}
\item The CD approached the FC for grant of a loan. The FC granted the CD a loan of 100 crores for setting up an SEZ project. The loan was secured by a mortgage and through a corporate guarantee furnished by the parent company, ACIL.
\item Once the CD defaulted, FC filed an application before the DRT to recover the outstanding loan amount.
\item A ``Debt Repayment and Settlement Agreement'' was consequently executed, in which the FC, the CD, and the corporate guarantor, CG (that is, ACIL), were parties.
\item Default was committed by the CD, and the FC invoked the corporate guarantee against CG, that is, ACIL. ACIL did not honour the guarantee, and an application under Section 7 was filed against ACIL.
\item The FC filed a claim for 650 crore, out of which 350 crore was admitted by the IRP. The IRP subsequently reassessed the claim at 241 crore, inclusive of the principal of 100 crores.
\item The resolution plan in the CIRP of the CG, that is, ACIL, was approved, and 38 crores were paid to the FC against the admitted claim of 247 crores in full and final settlement of its dues and demands.
\item Later, the FC filed an application under Section 7 against the CD. The FC filed a claim of 1,428 crore, which is the balance amount payable to the FC under the loan facility of 100 crores. The NCLT admitted the application under Section 7.
\item Aggrieved by the order, the appellant preferred an appeal before NCLAT. Both appeals were dismissed, and hence appeal to the SC.
\end{itemize}

The appellant is the SRA of the CG.

\section*{Appellant's Arguments}

Upon receipt of 38 crores from the guarantor, the debt repayable to the FC has been discharged. The FC is now estopped from enforcing the remaining part from the CD in view of Section 63 read with Section 41 of the Indian Contract Act.
\end{document}
